\chapter{Citations and links} \label{ch3}

References are managed with biblatex and biber, in the author--year style:
citations show the surname of the authors and the year of publication, and
the bibliography at the end of the thesis is sorted alphabetically.

\section{The bibliography file}

References are stored in \texttt{bibliography.bib}, one entry per work, each
identified by a key:
\begin{Verbatim}
@article{breiman2001random,
  author  = {Breiman, Leo},
  title   = {Random Forests},
  journal = {Machine Learning},
  volume  = {45},
  number  = {1},
  pages   = {5--32},
  year    = {2001},
  doi     = {10.1023/A:1010933404324}
}
\end{Verbatim}
Write author names as \texttt{Surname, Name} and separate them with
\texttt{and}. Export services such as Google Scholar, DBLP or the publisher's
web site provide entries in this format, but check them: they often contain
mistakes or miss the DOI. Only the works actually cited in the text are
listed in the bibliography, which is printed by \cs{printthesisbibliography}
at the end of \texttt{main.tex}.

\section{Citing}

Citations use the familiar commands of the natbib package. The two most
common forms are the parenthetical one, \cs{citep}, and the textual one,
\cs{citet}, where the author names are part of the sentence:
\begin{Verbatim}
Random forests \citep{breiman2001random} combine ...
As shown by \citet{breiman2001random}, ...
\end{Verbatim}
which produce: ``Random forests \citep{breiman2001random} combine \dots''
and ``As shown by \citet{breiman2001random}, \dots''. Table~\ref{tab:cite}
lists the other available commands.

\begin{table}[htbp]
  \centering
  \caption{Citation commands and their output.}
  \label{tab:cite}
  \begin{tabular}{ll}
    \toprule
    Command & Output \\
    \midrule
    \verb|\citep{cortes1995support}|            & \citep{cortes1995support} \\
    \verb|\citet{cortes1995support}|            & \citet{cortes1995support} \\
    \verb|\citep[p.~280]{cortes1995support}|    & \citep[p.~280]{cortes1995support} \\
    \verb|\citep[see][]{vapnik1995nature}|      & \citep[see][]{vapnik1995nature} \\
    \verb|\citep{bishop2006pattern,|            & \citep{bishop2006pattern,hastie2009elements} \\
    \verb|       hastie2009elements}|           & \\
    \verb|\citealp{pedregosa2011scikit}|        & \citealp{pedregosa2011scikit} \\
    \verb|\citeauthor{goodfellow2016deep}|      & \citeauthor{goodfellow2016deep} \\
    \verb|\citeyear{goodfellow2016deep}|        & \citeyear{goodfellow2016deep} \\
    \bottomrule
  \end{tabular}
\end{table}

Works with more than two authors are cited with the first author followed
by ``et al.'', as in \citet{hastie2009elements} or
\citet{pedregosa2011scikit}; all the authors are listed in the
bibliography.

\section{DOIs and URLs}

When an entry has a \texttt{doi} field, the bibliography shows it as a
\texttt{https://doi.org/\dots} link; otherwise, the \texttt{url} field is
shown, if present (see the entry of \citealp{goodfellow2016deep}). The link
closes the entry, right-aligned and never broken across lines. Prefer the
DOI to the URL whenever it exists: it is a permanent identifier, while web
addresses change over time.

Web addresses in the text are written with \cs{url}, which breaks them only
after characters such as slashes and dots. The address must be written
without spaces or line breaks:
\begin{Verbatim}
documented at \url{https://scikit-learn.org/stable/modules/...}
\end{Verbatim}
For instance, the documentation of the random forest classifier provided
by scikit-learn is available at
\url{https://scikit-learn.org/stable/modules/generated/sklearn.ensemble.RandomForestClassifier.html}.

\section{Links} \label{sec:links}

Citations, cross-references, URLs and DOIs are all hyperlinks, coloured in
the blue of the SLIM logo. This dark blue is barely distinguishable from
black on paper. With the option \texttt{pdfa=false}, links are blue when the
PDF is viewed on screen and black when it is printed, so that the printed
thesis has a uniform appearance; this is not allowed in a PDF/A file, and a
few PDF viewers that ignore the distinction may show the links in black or
print them in blue. To create a link with arbitrary text, use \cs{href}:
\begin{Verbatim}
the \href{https://slim.di.unimi.it}{SLIM laboratory}
\end{Verbatim}
which gives: the \href{https://slim.di.unimi.it}{SLIM laboratory}.
